\section{Full support for Gaussian prime conductors}
\label{sec:gaussian}

For a singular modulus $j$ of degree $h$, define $s(j)$ as the minimum
number of nonzero coefficients in a nonzero polynomial $P\in\mathbb Q[X]$
with $P(j)=0$. The minimal polynomial gives $s(j)\le h+1$. Equality
means that \emph{every} $h$ distinct nonnegative powers of $j$ form a
basis of $\mathbb Q(j)$. There is no bound on the exponents here.

\medskip\noindent\textbf{Theorem.}
For a rational prime $\ell$, let $j_\ell$ have endomorphism order
$\mathbb Z+\ell\mathbb Z[i]$, of discriminant $-4\ell^2$. The set of
primes with $s(j_\ell)=[\mathbb Q(j_\ell):\mathbb Q]+1$ has lower
relative natural density greater than $0.999818260107372245$.
The proof is independent of RH and GRH. The exact density of the
full-support set remains open.

\paragraph{The earlier problem.}
The ProveIt report \emph{Five-term equations for singular moduli},
in its sections on prime conductors and further questions, derives a
lower density greater than $0.6133$ and asks for the true density.
Its local calculation puts all conjugates in one residue class modulo
$7$. We use good integral models of the curves at $3$ and $7$ and an
explicit isogeny calculation at $2$ to obtain three multiplicative disks. Then we use all factorial comparisons across a single
cut. The finite density computations below concern the resulting
sufficient conditions. The
\href{https://github.com/VladimirReshetnikov/ProveIt/blob/6fe2c1bd450aa7f19eab17fb9a6e76170b4fce00/docs/incoming/grh_sparse_singular_moduli.zip}{prior archive}
contains \texttt{local.tex}, \texttt{prime-conductors.tex} and
\texttt{closing.tex}. For the original trinomial problem see
\href{https://arxiv.org/abs/2003.06547}{Bilu--Luca--Pizarro-Madariaga}.

\paragraph{The common disks at odd primes.}
We may discard $\ell=2,3,7$. Write $\sigma=\chi_{-4}(\ell)=\pm1$.
The order class number formula gives $h=(\ell-\sigma)/2$, since
$\mathbb Z[i]$ has class number one and the unit index is two.
Every conjugate is the invariant of a cyclic degree-$\ell$ quotient
of $E_0:y^2=x^3-x$. Indeed its CM lattice corresponds to an invertible
ideal $I$ of the order. Passing from $I$ to $I\mathbb Z[i]$ gives index
$\ell$, as follows locally from invertibility. The extended lattice
is principal, so the resulting curve has invariant $1728$. Taking
the dual isogeny gives the desired quotient. For background on this
lattice description see
\href{https://arxiv.org/abs/1208.5370}{Sutherland, \emph{Isogeny volcanoes}}.

At $p=3,7$, the curve $E_0$ has good reduction. In each case a finite
unramified extension contains all its $\ell$-torsion, since the torsion
of order prime to $p$ is finite \'etale on the good model. Over its
integer ring $\mathcal O$, quotienting by each cyclic subgroup gives
an elliptic curve with good reduction. Its reduction is supersingular
because it is isogenous to the supersingular reduction of $E_0$.
The latter assertion also follows from the Hasse invariant computation
below applied to $E_0$. The good-reduction facts can equivalently be
deduced from the N\'eron--Ogg--Shafarevich criterion.

At $p=7$ use a good integral short equation $y^2=x^3+Ax+B$.
The Hasse invariant is the coefficient of $x^6$ in $(x^3+Ax+B)^3$,
which is $3B$. Thus $7\mid B$. The denominator $4A^3+27B^2$ is a
unit, and
\[
 \frac{j}{1728}=\frac{4A^3}{4A^3+27B^2}
                 =1-\frac{27B^2}{4A^3+27B^2}\in 1+7^2\mathcal O .
\]
At $p=3$, complete the square in a good integral equation to get
$y^2=x^3+a_2x^2+a_4x+a_6$. The Hasse invariant is the coefficient of
$x^2$ in this cubic, so $3\mid a_2$. Translating the $x$ coordinate
by $-a_2/3$ gives an integral short equation with unit discriminant.
The same formula now gives $j/1728\in 1+3^3\mathcal O$. Division by
$1728$ is taken in the local field at $3$.

\paragraph{A common multiplicative disk at 2.}
Let $v(2)=1$. For any two cyclic order-$\ell$ subgroups $H,J$ of
$E_0$, with $\ell$ odd, we claim
\[
 v\big(j(E_0/H)/j(E_0/J)-1\big)\ge4.                 \tag{G0}
\]
All valuations may be taken in one finite extension of $\mathbb Q_2$.
The integral invariant $1728$ gives $E_0$ potentially good reduction.
Take a field of good reduction containing $E_0[3\ell]$. Write
$x=r+u^2X$ for the change from the original short equation to a good
integral Weierstrass equation with coordinate $X$. Its discriminant
is a unit, while the original discriminant is $64$, so $v(u^2)=1$.
Every nonzero odd-torsion point under consideration has integral $X$:
reduction is injective on prime-to-$2$ torsion, so its reduction is
an affine point of the good model. In particular for a fixed
nonzero $3$-torsion point $T$ and every nonzero $Q\in H\cup J$ we
have $v(x_Q-x_T)\ge1$.

The $3$-division polynomial gives $3x_T^4-6x_T^2-1=0$.
Its roots are units. For $y=5x_T^2-3$ this yields
$3y^2-12y-88=0$. The Newton polygon has its only slope $-3/2$,
so $v(y)=3/2$. Since $x_Q=x_T+2z$ with $z$ integral,
$v(x_Q^2-x_T^2)\ge2$. For $g(x)=x(5x^2-3)$ it follows that
\[
 v(g(x_T))=3/2,\qquad v(g(x_Q)-g(x_T))\ge2.
\]
V\'elu's quotient model is $Y^2=X^3+A_H X+B_H$, where
\[
 A_H=-1-10\!\sum_{Q\in(H\setminus\{0\})/\{\pm1\}}(3x_Q^2-1),
 \qquad
 B_H=-14\!\sum_{Q\in(H\setminus\{0\})/\{\pm1\}}g(x_Q).
\]
These are the short-equation specialization of
\href{https://ddd.uab.cat/pub/pubmat/pubmat_a2007vEXTRA/pubmat_a2007vextrap147.pdf}{Miret--Moreno--Rio,
\emph{Generalization of V\'elu's formulae for isogenies between elliptic curves},
pp.~149--151, especially Theorem~1}.
In particular $A_H$ is a unit congruent to $-1$ modulo $2$ and
$v(B_H)\ge5/2$. The $H$ and $J$ sums have the same number
$(\ell-1)/2$ of terms, all congruent to $g(x_T)$ modulo $4$ in
the second sum. Thus $v(B_H-B_J)\ge3$ and
$v(B_H+B_J)=v(B_H-B_J+2B_J)\ge3$.
This gives $v(B_H^2-B_J^2)\ge6$.

Put $s_H=27B_H^2/(4A_H^3)$. Then $v(s_H)\ge3$ and
$v(s_H-s_J)\ge4$: replacing $B_H^2$ by $B_J^2$ costs at least four
in valuation after division by $4$, while
$v(A_H^{-3}-A_J^{-3})\ge1$. The identity
$j(E_0/H)=1728/(1+s_H)$ proves (G0).
Fixing any one of the quotient invariants $j_*$ therefore puts all
$h$ conjugates in $j_*(1+16\mathcal O)$, where $\mathcal O$ is the
valuation ring. The center is allowed to be algebraic. The recurrence
and the root count below apply over this complete finite extension
with the same normalized valuation.

The depth four is sharp when the disk assertion ranges over all odd
prime conductors. For $\ell=3$ let $Y=j/1728-1$. Substitution of
$A=9-30x_T^2$, $B=-14x_T(5x_T^2-3)$ into the formula for $j$
and elimination by $3x_T^4-6x_T^2-1=0$ give the quadratic
$27Y^2-2399040Y-18593344$. Its discriminant has $2$-adic valuation
eight. The two roots differ by an element of valuation four, and
both $1+Y$ are units. The corresponding pair thus attains equality
in (G0). The archive reproduces the elimination and also checks
the corresponding degree-five isogenies in characteristic zero.


\paragraph{The factorial comparison.}
Use $(p,b)=(2,4),(3,3)$ or $(7,2)$ and $w(n)=bn-v_p(n!)$, with $w(0)=0$.
The sufficient condition at $p$ is
\[
       \max_{0\le u<h}w(u)<\min_{v\ge h}w(v).                 \tag{G1}
\]
To prove this, suppose that a polynomial with $k\le h$ terms and
distinct exponents $e_j$ vanishes at a conjugate. It vanishes at all
$h$ conjugates. Let $u_0$ be the common multiplicative center given by these disks,
namely $1728$ at $3,7$ and a chosen conjugate at $2$. Compose with
$X=u_0\exp(p^bY)$ on the closed unit disk of a finite local field
containing the center and all conjugates. The logarithm gives a distinct point
$Y$ for each root. The coefficient at degree $n$ is $p^{bn}A_n/n!$,
where $A_n=\sum_j c_ju_0^{e_j}e_j^n$. A Vandermonde determinant
gives $M=\max_{n<k}|A_n|_p>0$. The monic integral recurrence with
characteristic polynomial $\prod_j(T-e_j)$ implies $|A_n|_p\le M$
for all $n$. Under (G1), all coefficients from $n=h$ onward have
smaller norm than the coefficient at an index below $k$ attaining
$M$. The series converges, since $b>1/(p-1)$.
Strassmann's theorem bounds the number of its roots by $h-1$, a
contradiction. This is the disk argument of the prior report applied
to the deeper common disks. For context on fewnomial root counts see
\href{https://arxiv.org/abs/1008.4808}{Avenda\~no--Krick}.

\paragraph{Counting the prime conductors.}
Taking differences $w(v)-w(u)$ shows that (G1) is equivalent to
\[
        \sum_{t\in I}v_p(t)<b|I|
        \quad\text{for every integer interval }I\subset\mathbb Z_{>0}
        \text{ containing }h.                                \tag{G2}
\]
For lengths $r\le R$, omitting the finitely many $h<R$, the test
depends only on $h\bmod p^{bR}$. Fix $\sigma$. For an odd $p$, the admissible classes of $h$ modulo $p^a$
are those with $p\nmid 2h+\sigma$. For $p=2$, they are exactly the
even classes: $\ell=2h+\sigma\equiv\sigma\pmod4$,
and all such classes correspond to invertible classes of $\ell$
modulo $2^{a+1}$ with that sign. In all three cases each admissible class
has conditional density $1/\varphi(p^a)$, by the prime number theorem
in arithmetic progressions. Reflection $h\mapsto-h$ and reversal of
intervals show that the finite density $d_{p,R}$ is the same for
either $\sigma$.

Here is a uniform bound passing from finite lengths to the infinite
test. If a length-$r$ interval violates (G2) and its greatest
valuation is $m$, counting multiples of $p^a$ for $1\le a\le m$
bounds the sum of valuations by $r/(p-1)+m$. Consequently some
$h+j$, with $|j|\le r-1$, is divisible by
\[
      q_r=p^{\lceil (b-1/(p-1))r\rceil}.
\]
For fixed $r$ the conditional upper density of this event is at most
$u_r=(2r-1)/\varphi(q_r)$. To justify summing the bound, write
$X$ for the prime cutoff. For $q_r\le\sqrt X$ the Brun--Titchmarsh
inequality bounds the possible primes in each relevant invertible residue
class by $O(X/(\varphi(q_r)\log X))$, with an absolute constant.
The modulus is $q_r$ for odd $p$ after dropping the condition modulo
$4$, and $2q_r$ for $p=2$, because $\ell=2h+\sigma$. In the latter
case $2q_r\le 2\sqrt X$ still gives $\log(X/(2q_r))\gg\log X$. Since only $\ell=p$
can occupy a noninvertible class, dropping that prime and dividing
by the count of primes of each sign gives a uniform tail bounded by
an absolute multiple of $\sum_{r>T}u_r$. For $q_r>\sqrt X$, a positive
multiple of $p^{\lfloor\log_p\sqrt X\rfloor}$ occurs among the $h+j$.
The inequality $(b-1/(p-1))r\le\log_p(h+r)$ forces $r=O(\log X)$:
when $r>h$ the inequality fails since $h+r<2r$ and
$2r\le 2^r<p^{(b-1/(p-1))r}$. Otherwise $h+r\ll X$.
There are $O(\sqrt X\log X)$ possibilities for a multiple and an
offset, negligible compared with the number of primes. Taking limits
at fixed $T$ and then $T\to\infty$ proves existence of the test
density $d_p$ and
\[
         0\le d_{p,R}-d_p\le\sum_{r>R}u_r.                  \tag{G3}
\]
The sieve inequality used here is due to
\href{https://doi.org/10.1112/S0025579300004708}{Montgomery--Vaughan}.

\paragraph{Exact certificate.}
The verifier starts with the $p-1$ admissible classes modulo $p$.
In a class modulo $p^a$, it evaluates all relevant valuations truncated
at $a$. If any interval sum already reaches $br$ it discards the
class. If no valuation is truncated it accepts the class. Otherwise
it visits all $p$ lifts. The process terminates by depth $bR$, since
a truncated term at that depth already violates a length at most $R$
containing that term and the center. To bound (G3), group its terms
in successive blocks of length $p-1$. For $r\ge R+1$,
\[
 \frac{u_{r+p-1}}{u_r}=
 \frac{2r+2p-3}{2r-1}\,p^{-(b(p-1)-1)}
\]
is decreasing. A geometric sum bounds the whole tail using its first
block. All operations in \texttt{check\_density.py} use rational
arithmetic for the bounds.
\[
\begin{array}{c|c|c|c}
p & R & d_{p,R} & \text{tail upper bound}\\ \hline
2&25&65963232382760376856543/75557863725914323419136 &3.879\cdot10^{-22}\\
3&16&68736524642595367189/72945992754341572806 &1.694\cdot10^{-19}\\
7&10&9528160153068835/9770481587462694 &5.089\cdot10^{-17}
\end{array}
\]
At fixed sign, the three finite congruence tests are independent by the
Chinese remainder theorem and the prime number theorem in progressions.
Passing to the limits using (G3), the density of passing at least one test
is $d=1-(1-d_2)(1-d_3)(1-d_7)$. It is the same for both signs and satisfies
\[
 0.9998182601073722452943 < d < 0.9998182601073722456678 .
\]
Every prime counted by $d$ gives full support. Determining what occurs
for the remaining conductors would require information beyond these
three disk tests.

\paragraph{Individual conductors.}
For $\ell=197$ the degree is $98$. The one-term interval $\{98\}$
fails (G2) at $7$, while it passes at $3$. At $3$ every interval of
length $r\ge2$ through $98$ has valuation sum at most
$r/2+\log_3(98+r)<3r$. The last inequality holds at $r=2$ and
persists since $3^{5r/2}$ grows faster than $98+r$ at every step.
Thus any $98$ distinct nonnegative powers of $j_{197}$ form a basis.
For $\ell=107$ the degree is $54$. Here $3^3\mid54$ rules out (G2)
at $3$, while $7\nmid54$. For $r\ge2$ one has
$r/6+\log_7(54+r)<2r$, again by checking $r=2$ and increasing $r$.
Every $54$ distinct powers of $j_{107}$ therefore form a basis.

Conductor $15877$ exhibits the contribution at $2$. It is prime
(by trial division through $126$) and $h=7938=2\cdot 3^4\cdot7^2$.
Both odd-prime tests fail on the singleton $\{h\}$. The largest
sums of $2$-adic valuations of intervals through $7938$ of lengths
$1,2,3,4$ are respectively $1,1,9,9$, each smaller than four times
the corresponding length. For every $r\ge5$, the sum is at most
$r+\log_2(7938+r)<4r$, since $7938+5<2^{15}$ and the inequality
persists on increasing $r$. Thus any $7938$ distinct nonnegative
powers of this singular modulus form a basis.
